\section{What GPUs can and cannot accelerate}
\label{sec:gpu}

GPUs are good at large batches of similar operations and less well suited to
irregular, sequential control flow. SAT offers both. Batched assignment
evaluation, independent local searches, message updates and some simplification
kernels expose parallel work. CDCL's trail, propagation dependencies, changing
clause database and divergent search paths make a straightforward wholesale
port difficult. The relevant question is therefore which computation is
accelerated and what fraction of total solving cost it represents.

\subsection{Published evidence}
\textbf{GPU inprocessing inside a complete solver.}
ParaFROST accelerates simplification and inprocessing and supports proof
logging \citep{osama2024}. Its experiments use an Intel Xeon Gold 6130,
96 GB RAM and a Titan RTX, with a 3600-second solving limit.
The CPU comparisons use a different Xeon E5-2630 machine, so the hardware
comparison is not fully matched.
The GPU version improves on its related CPU engine, but the paper also
reports that Kissat is faster overall on the evaluated collection.
This supports GPU acceleration of selected components, not an overall
advantage against every strong CPU solver. Proof production and checking
cost must be included when certification is the objective.

\textbf{Continuous local search.}
FastFourierSAT maps gradient computation for symmetric Boolean constraints
to GPU-friendly operations \citep{cen2025}. It is particularly relevant to
hybrid formulas containing native cardinality and XOR constraints.
The paper reports large kernel gains, but its ordinary CNF experiment is
cautionary: in a 1000-second test on satisfiable SAT Competition 2023
instances, its GPU variant solves 13 and its CPU variant 5. The corresponding
competition solvers solve those easy instances in roughly $0.04$--$1.07$
seconds. The GPU setup includes an RTX 3080 Ti. Thus a large acceleration
within one continuous-search method does not establish competitiveness with
CDCL on general CNF. The method is incomplete and cannot certify UNSAT.

\subsection{Recent preprints and comparison limits}
\textbf{TurboSAT.}
The November 2025 preprint combines GPU continuous search with CPU CDCL
\citep{dai2025}. On its 179-instance SAT Competition 2024 collection it
reports 165 solved within 5000 seconds, versus 143 for the chosen CaDiCaL
1.9.5 baseline. Experiments use a DGX GB200 system with up to two B200 GPUs
and 141 CPU threads. The method retains an unrestricted CDCL worker for
completeness. This is promising hybrid-system evidence, but it is not an
isolated GPU speedup: the CPU portfolio, hardware resources and configurations
also differ. These benchmarks are not a controlled random-$k$-SAT size study.

\textbf{Accelerated Fourier SAT.}
The June 2026 AFSAT preprint studies GPU implementation and throughput,
including hard random $3$-SAT and scaling across V100 GPUs
\citep{christopher2026}. It helps establish that this style of search can
benefit from parallel hardware. Throughput, however, must be combined with
the probability of finding a fully satisfying assignment to obtain a
time-to-solution comparison. Continuous objective values and almost-satisfying
outputs are insufficient; the benchmark must enforce zero violated clauses.

\begin{table}[htbp]
\centering\small
\begin{tabularx}{\linewidth}{@{}L{.27\linewidth}Y Y@{}}
\toprule
Approach & Evidence of benefit & Main limitation \\
\midrule
ParaFROST & GPU simplification within complete solving; published & Improvement over related CPU engine does not imply beating Kissat. \\
FastFourierSAT & Parallel gradients, especially hybrid constraints; published & Kernel gains and native-constraint advantages need not transfer to random CNF. \\
TurboSAT & More solves in a large hybrid system; preprint & Many CPU workers and GPUs confound a GPU-only comparison. \\
AFSAT & GPU throughput and random-$3$-SAT experiments; preprint & Throughput is not a complete-decision or UNSAT guarantee. \\
\bottomrule
\end{tabularx}
\caption{Four concrete GPU approaches. Publication status and benchmark
limitations are part of the comparison, not interchangeable speedup claims.}
\end{table}

\subsection{Recommendation for this project}
Establish a CPU baseline first. A GPU comparison is most informative for
large satisfiable random instances, many independent instances or restarts,
and batched message-passing or continuous-search experiments. Keep formula
generation, preprocessing, host--device transfer, compilation/warm-up,
assignment reconstruction and verification in the accounting. Compare
matched hardware budgets as well as raw wall-clock time; report CPU core
counts and GPU counts explicitly.

For near-threshold UNSAT certification, the checked literature provides no
basis for assuming that buying a GPU removes the principal difficulty.
GPUs may help specific kernels or portfolios, but a reduction in kernel
time cannot by itself overcome a search or proof-size bottleneck. A useful
reported speedup is the time to a verified answer on the same instances at
the same success target, not only flips per second or gradient evaluations
per second.
